% Part: model-theory

\documentclass[../../include/open-logic-part]{subfiles}

\begin{document}

\olpart{mod}{నమూనా సిద్ధాంతం}

\begin{editorial}
  నమూనా సిద్ధాంతానికి సంబంధించిన పాఠ్యసామగ్రి అసంపూర్ణంగానూ,
  ప్రయోగాత్మకంగానూ ఉంది. ప్రస్తుతం ఇది ఆల్డో ఆంటొనెల్లి నమూనా
  సిద్ధాంత గమనికలకు చేసిన అనుసరణ మాత్రమే; మొదటిస్థాయి తర్క భాగంలో
  ఇప్పటికే చర్చించిన సిద్ధాంతాలు, సంపూర్ణత, సంహతత్వం వంటి అంశాలను
  ఇందులో మినహాయించారు. పాఠ్యపుస్తకంలో ఉపయోగపడాలంటే దీనికి ఇంకా ఎంతో
  పరిచయం, ప్రేరణ, వివరణతోపాటు అభ్యాసాలు అవసరం. ఆండీ అరానా ఈ భాగంపై
  ప్రత్యేకంగా పనిచేయాలని యోచిస్తున్నారు (\gitissue{65}).
  \sourcecorrection{OLTEMODBAS-001}{ఆంగ్ల మూలంలోని \texttt{is at planning to work} అనే వ్యాకరణదోషాన్ని, సందర్భం కోరే \texttt{is planning to work}గా చదివి అనువదించాం; వ్యక్తి పేరు, సమస్య సూచిక, ఉద్దేశం మారలేదు.}
\end{editorial}

\olimport[basics]{basics}

\olimport[models-of-arithmetic]{models-of-arithmetic}

\olimport[interpolation]{interpolation}

\olimport[lindstrom]{lindstrom}

\OLEndPartHook

\end{document}
